% !TeX spellcheck = en_US
\documentclass[french]{yLectureNote}

\title{Électromagnétisme 2 PS}
\subtitle{Physique}
\author{Paulhenry Saux}
\date{\today}
\yLanguage{Français}
%lionels.calmels@cemes.fr
\professor{M.Brut}
\usepackage{graphicx}%----pour mettre des images
\usepackage[utf8]{inputenc}%---encodage
\usepackage{geometry}%---pour modifier les tailles et mettre a4paper
%\usepackage{awesomebox}%---pour les boites d'exercices, de pbq et de croquis ---d\'esactiv\'e pour les TP de PC
\usepackage{tikz}%---pour deiffner + d\'ependance de chemfig
% \usepackage{tabularx}%---pour dimensionner automatiquement les tableaux avec variable X
\usepackage{awesomebox}%---Pour les boites info, danger et autres
\usepackage{menukeys}%---Pour deiffner les touches de Calculatrice
\usepackage{fancyhdr}%---pour les en-t\^ete personnalis\'ees
\usepackage{blindtext}%---pour les liens
\usepackage{hyperref}%---pour les liens (\`a mettre en dernier)
\usepackage{caption}%---pour la francisation de la l\'egende table vers Tableau
\usepackage{pifont}
\usepackage{array}%---pour les tableaux
\usepackage{lipsum}
\usepackage{yFlatTable}
\usepackage{multicol}
\newcommand{\Lim}[1]{\lim\limits_{\substack{#1}}\:}
\renewcommand{\vec}{\overrightarrow}
\newcommand{\N}[0]{\mathbb{N}}
\newcommand{\dd}{\mathrm{d}}
\newcommand{\norm}[1]{||\vec{#1}||}
\newcommand{\fo}{\psi(\vec{r},t)}
\newcommand{\foe}{\psi(\vec{r},t)\*}
\newcommand{\HH}{\hat{H}}
\newcommand{\hb}{\hbar}
\newcommand{\lap}{\nabla^2}
\newcommand{\lapcc}{\frac{\partial^2 }{\partial x^2}+\frac{\partial^2 }{\partial y^2}+\frac{\partial^2 }{\partial z^2}}
\newcommand{\E}{\vec{E}(M)}
\newcommand{\B}{\vec{B}(M)}
\newcommand{\V}{V(M)}
\newcommand{\er}{\vec{e_{\rho}}}
\newcommand{\ep}{\vec{e_{\varphi}}}
\newcommand{\pip}{\pi^+}
\newcommand{\pim}{\pi^-}
\newcommand{\mv}{\mathcal{V}}
\DeclareMathOperator{\Div}{Div}
\newcommand{\rot}{\vec{\mathrm{rot}}}
\newcommand{\grad}{\vec{\mathrm{grad}}}
\setcounter{chapter}{1}
\begin{document}
\chapter{Induction}
\section{Lois de l'induction}
\subsection{Force électromotrice}
\begin{definition}[Force électromotrice (fem)]
    Toute variation de flux magnétique à travers un circuit provoque l'apparition d'une fem. Le circuit est alors équivalent à
    une source de tension notée $e$. Si le circuit est ouvert, la tension vaut $e$, dans le cas contraire, elle induit l'apparition d'un courant induit noté $i$.

    $$e(t) = \int \E_m\cdot \vec{\dd l}$$ où $\E_m$ est le champ électromoteur
\end{definition}
En utilisant la force de Lorentz et l'équation de Maxwell-Faraday, on décompose $\vec{E}_m$ selon l'origine physique du mouvement des charges :
\begin{theorem}[Décomposition de la fem]
    $$e(t) = \oint_C \left( \frac{\partial \vec{B}}{\partial t} + \vec{v} \wedge \vec{B} \right) \cdot \vec{dl}$$
\end{theorem}
\begin{itemize}
    \item \textbf{Terme de Neumann :} Lié à la variation temporelle du champ magnétique ($\vec{\text{rot}} \vec{E}_{induit} = -\frac{\partial \vec{B}}{\partial t}$). C'est l'induction "statique".
    \item \textbf{Terme de Lorentz :} Lié au mouvement du circuit à la vitesse $\vec{v}$. C'est l'induction "cinématique".
\end{itemize}
\subsection{Lois de Faraday}
\begin{theorem}[Loi de Faraday]
    $$e(t) = -\frac{\dd \phi}{\dd t}$$ avec $\phi = \iint_S \B\cdot \vec{\dd S}$.
\end{theorem}
\warningInfo{Validité de la loi de Faraday}{On peut croire qu'elle est valable uniquement quand le circuit est immobile. Cependant, si le circuit bouge, la surface S sur laquelle est calculé l'intégrale dépend du temps, ce qui est pris en compte lors de la dérivation.}
\begin{theorem}[Loi de Lenz]
    Les phénomènes d'induction s'opposent par leur effets aux causes qui leur ont donné naissance. On en déduit que le sens du
    courant induit est tel que le flux de son propre champ magnétique s'oppose à la variation de flux inducteur.
\end{theorem}
\begin{theorem}[Équation de Maxwell Faraday]
    $$\rot \E_{induit} = -\frac{\partial \B}{\partial t}$$
\end{theorem}
\criticalInfo{Différences de champ}{
    \begin{itemize}
    \item \textbf{Le champ induit $\vec{E}_i$ :} Champ électrique d'origine magnétique, lié à la variation temporelle de $\vec{B}$ (induction de Neumann). Il vérifie $\vec{\text{rot}} \vec{E}_i = -\frac{\partial \vec{B}}{\partial t}$.
    \item \textbf{Le champ électromoteur de Lorentz $\vec{v} \wedge \vec{B}$ :} Contribution liée au mouvement du circuit dans un champ magnétique (induction de Lorentz).
    \item \textbf{Le champ électromoteur total $\vec{E}_m$ :} Somme des deux contributions précédentes. Sa circulation le long du circuit donne la f.e.m. totale :
    $$e(t) = \oint_C \vec{E}_m \cdot \vec{dl} = \oint_C (\vec{E}_i + \vec{v} \wedge \vec{B}) \cdot \vec{dl}$$
\end{itemize}
}
\subsection{Distinction entre Dérivée Partielle et Dérivée Totale}
\subsubsection{L'approche locale : On dérive $\vec{B}$ puis on intègre}
C'est ce que l'on fait pour calculer la contribution de \textbf{Neumann} uniquement. On regarde comment le champ change en chaque point d'une surface que l'on considère comme "gelée" à l'instant $t$ :
$$\text{f.e.m. de Neumann} = \iint_{S} \left( -\frac{\partial \vec{B}}{\partial t} \right) \cdot \vec{dS}$$
\begin{itemize}
    \item Ici, on utilise la \textbf{dérivée partielle} ($\partial$).
    \item On ne s'occupe pas du mouvement du contour. 
    \item Cette intégrale est égale à la circulation du champ induit : $\oint_C \vec{E}_i \cdot \vec{dl}$.
\end{itemize}

\subsubsection{L'approche globale (Faraday) : On intègre puis on dérive le tout}
C'est la méthode la plus complète. On calcule d'abord le flux total $\phi(t)$ à travers la surface $S(t)$ telle qu'elle est à l'instant $t$ (en tenant compte de sa position et de sa forme), puis on dérive le résultat par rapport au temps :
$$e(t) = -\frac{d}{dt} \left( \iint_{S(t)} \vec{B} \cdot \vec{dS} \right)$$
\begin{itemize}
    \item Ici, on utilise la \textbf{dérivée totale} ($d/dt$).
    \item Cette opération mathématique "capture" automatiquement deux effets :
    \begin{enumerate}
        \item La variation de $\vec{B}$ (le terme en $\partial \vec{B} / \partial t$ ci-dessus).
        \item Le mouvement de la surface (qui correspond physiquement au terme de Lorentz $\vec{v} \wedge \vec{B}$).
    \end{enumerate}
\end{itemize}
\section{Auto-Induction et inductance d'une bobine}
\subsection{Phénomène d'auto-Induction}
On considère une spire avec une inensité variable. On observe le champ crée par l'inetsnité mais il y a aussi le 
champ induit par la variation du courant. Cela crée un nouveau courant, qui va s'opposer au courant initial. C'est l'auto-induction.

\begin{definition}[Auto-induction]
    La spire parcourue par un courant crée un champ B qui crée un courant induit dans cette spire.
    Par la loi de Lorenz, le courant induit s'oppose au courant initial.
\end{definition}
\subsection{Induction propre et inductance}
\begin{definition}[Flux propre]
    Flux crée par la spire à l'intérieur d'elle-m\^eme, dépend de B, qui dépend des caractéristiques
    de la spire et de l'intensité du courant.
\end{definition}
On peut donc écrire $\phi = L i(t)$, avec L qui est une constante en henry. 

On a donc $$e= -\frac{\dd \phi}{\dd t} = -L\frac{\dd i}{\dd t}$$

\end{document}
